基因编辑效率的机器学习预测已取得大量进展，但绝大多数报告只给出单一数据集上的精度，既未检验该精度能否在\textbf{独立数据集}上复现，也未检验模型归因所指认的序列位置能否被\textbf{独立第三方模型}证实。更关键的是，若训练与测试之间存在同源序列，跨条件性能会被系统性高估，而这类问题在只报告精度时不可见。本文以"多数据集复现 $\rightarrow$ 多模型外部验证 $\rightarrow$ 可解释分析 $\rightarrow$ 反事实扰动 $\rightarrow$ 生物学假设"为证据链，在一个经过完整质量审计、序列级无泄漏的计算平台上回答三个问题：(1) 序列可预测性在不同数据集间是否复现？(2) 多模型归因指认的序列位置能否被独立外部模型证实？(3) 该位置能否转化为可证伪的最小扰动假设？

本文使用三个相互独立的公开数据集——DeepCRISPR（4 个细胞系、16\,749 条 23\,nt 靶序列、1\,344 次受控训练运行）、Hiranniramol（1\,309 条）与 Labuhn（417 条）——在同一套特征构造、同一批 7 个模型配置（Linear、XGBoost、MLP、双分支 CNN 卷积核 $k\in\{3,5,7\}$、Transformer）与同一条身份类划分约束下训练，并引入独立第三方模型 CRISPRon 作为外部验证器。全文严格区分\textbf{六层证据}：预测证据（E1）、关联证据（E2）、归因证据（E3）、反事实证据（E4）、实验证据（E5）与因果证据（E6）；本文不含任何 E5 证据，因此不产生任何 E6 层面的断言。

主要结果如下。\textbf{第一，序列可预测性是数据集依赖的，不可假定普遍复现。} 在留出测试集上，DeepCRISPR 单细胞系划分的 7 个配置中位 $R^{2}$ 仅为 $0.067$--$0.120$；Hiranniramol 在完全相同的流程下达到 $0.345$--$0.489$（Transformer 最高，$R^{2}=0.489$，Pearson $=0.709$）；而 Labuhn 的 7 个配置全部为负（$-0.725$ 至 $-0.050$），模型预测坍缩为近常数。与官方划分无关的全数据 5 折交叉验证复核给出同向结论：Hiranniramol 的 CV $R^{2}=+0.389$ 远高于其标签置换零分布（$-0.467\pm0.038$，$n=200$），而 Labuhn 的 CV $R^{2}=-0.144$ 仅比自身零分布（$-0.287\pm0.050$）高 $0.14$。因此我们把"序列信号存在"限定为\textbf{数据集特异性结论}，而非普适结论。\textbf{第二，跨细胞系迁移失败。} 真实留一细胞系（LOCO）划分下，7 个配置的中位 $R^{2}$ 全部落在 $-0.027$ 至 $+0.010$，与均值基线不可区分。\textbf{第三，表观环境通道未提供可检出的增量预测信息。} CTCF、DNase、H3K4me3、RRBS 的主效应 $|\overline{\Delta R^{2}}|$ 为 $3.2\times10^{-4}$--$1.5\times10^{-3}$，比预设效应门 $0.01$ 低 1--2 个数量级，跨模型 bootstrap 区间全部跨 0，区组析因方差分析主效应 $p=0.975$--$0.9996$，四者均判为 Inconclusive。

\textbf{第四，也是本文证据链的收敛点：一个可由外部模型独立证实的位置效应。} 四个非模型类（XGBoost、CNN、MLP、Transformer）的归因峰值在 65.6\%--100\% 的上下文中落在 PAM 邻近窗口（17--20），其单模型平均谱峰值分别位于第 18、20、18、19 位；线性模型例外（21.9\%，峰值在第 1 位），因此把线性模型纳入等权平均后，跨模型平均谱峰值是第 1 位而非第 18 位（E3）。我们按\textbf{预先登记、不接触任何突变侧信息}的规则从 DeepCRISPR 测试集中选出 8 条第 18 位为 C 的代表性野生型序列，仅替换第 18 位一个碱基（C$\rightarrow$A），并用独立第三方模型 CRISPRon 评分（E4，外部模型验证）。结果显示：本文 7 个模型与 CRISPRon 的扰动响应在 8 条序列中的 \textbf{7 条上方向一致}（Pearson $r=0.629$，Spearman $\rho=0.476$），唯一分歧序列被保留而非剔除；CRISPRon 侧 7/8 条给出负向变化（平均 $\Delta=-12.43$）。进一步对全部 23 个位点做饱和 in-silico 突变（552 条替换记录，其中 48 条破坏 PAM 被排除），第 18 位的 C$\rightarrow$A 平均变化为 $-12.425$，而其余位置平均仅 $+0.144$；按平均 $|\Delta|$ 计第 18 位在 21 个可比位置中\textbf{排名第 1}，且逐序列看有 7/8 条排在第 1。与之独立的观测性关联（E2）方向一致但具细胞背景依赖：第 18 位 C 减 A 的实测平均效率差为 HCT116 $+0.090$、HeLa $+0.092$、HL60 $+0.027$、HEK293T $-0.015$。

综上，本文的贡献不是报告一个高精度预测器——事实上我们报告了该预测能力在不同数据集间的\textbf{复现失败}——而是在严格区分证据层级的前提下，把"序列可预测性有多可信"与"哪个序列位置值得实验检验"这两个问题分开回答：前者是数据集依赖的，后者在位置层级上获得了跨五类模型、跨独立外部模型与跨观测性关联的三方支持，因而可以转化为一个预登记可证伪的最小扰动假设（第 18 位单碱基替换）。本文不含任何湿实验数据，全部结论的最高证据层级为 E4，不构成因果断言。作为方法学附带的整改结果，我们还纠正了自身历史批次中跨细胞系同源序列污染带来的性能高估（\emph{mixed} 中位 $R^{2}$ 由 0.118 降至 0.073，LOCO 由 $+0.036$ 降至 $-0.008$），并如实报告该下降。
